\zad{1}{}{На выборах за кандидата А подано $a$ голосов, за Б~— $b$ голосов, $a>b$. Бюллетени перемешали и достают по одному, пока не достанут все, и каждый раз объявляют счёт. С какой вероятностью А после каждого бюллетеня был строго впереди? \par\vspace{3pt}К этой задаче мы ещё не раз вернёмся.}

\begin{samepage}
\zadp{2}{\circ}
\punkt{а}{}{Какие здесь элементарные исходы? Сколько их и чему равна вероятность каждого?}
\punkt{б}{}{Найдите ответ в задаче 1 при $a=3$, $b=2$ и при $a=4$, $b=2$.}
\punkt{в}{}{Как изобразить исход, чтобы с первого взгляда было видно, что А всё время был впереди?}
\end{samepage}\par\vspace{9pt}

\celoe{\tekst{\term{Путь}~— ломаная из звеньев $\nearrow$ и $\searrow$, \term{длина} пути~— число звеньев. Если начало пути не указано, он начинается в $(0,0)$. Вот десять путей из $(0,0)$ в $(8,2)$. Сколько всего таких путей?}\risunok{@puti10@}}

\risunok{@treugolnik@}\vspace{4pt}

\begin{samepage}
\zadp{3}{\circ}
\punkt{а}{}{В каких точках может кончаться путь Дика?}
\punkt{б}{}{Заполните треугольник Каталана. По какому правилу число в точке получается из соседних?}
\punkt{в}{}{Найдите ответ в задаче 1 при всех $a>b$, $a+b\leqslant 8$. Предположите, чему он равен в~общем~случае.}
\end{samepage}\par\vspace{9pt}

\celoe{\tekst{Число путей Дика в $(2n,0)$ называют \term{числом Каталана} и обозначают $C_n$. На рисунке~— все пути Дика~в~$(8,0)$.}\risunok{@dik14@}}

\begin{samepage}
\zadp{4}{\circ}
\punkt{а}{}{Скажем, что путь Дика впервые возвращается на ось в точке $(2k,0)$, если между началом и этой точкой он оси не касается. Сколько путей Дика в $(2n,0)$ впервые возвращаются на~ось в~точке~$(2k,0)$?}
\punkt{б}{}{Выразите $C_{n+1}$ через $C_0,\ldots,C_n$.}
\punkt{в}{}{Найдите ответ в задаче 1 при $a=n+1$, $b=n$.}
\end{samepage}\par\vspace{9pt}

\tekst{Дадим ещё три определения.}

\celoe{\tekst{$\bullet$~\term{Триангуляция} выпуклого многоугольника~— разрезание его непересекающимися диагоналями на треугольники. Множество триангуляций правильного $(n+2)$-угольника обозначим $\mathcal T_n$; триангуляции, переходящие друг в друга при повороте или симметрии, различаются. На рисунке~— все элементы $\mathcal T_3$.}\risunok{@triang5@}}

\celoe{\tekst{$\bullet$~\term{Двоичное дерево}~— дерево с выделенной вершиной (корнем), в котором у каждой вершины либо нет потомков, либо их ровно два, причём один из них назван левым, а другой~— правым. Вершины с потомками~— \term{развилки} (крупные точки). Множество двоичных деревьев с $n$ развилками обозначим $\mathcal D_n$. На рисунке~— все элементы $\mathcal D_3$.}\risunok{@derevya5@}}

\sboku{8.2cm}{@tupik@}{\tekst{$\bullet$~Рассмотрим железнодорожный тупик, справа от которого стоят вагоны, пронумерованные числами от 1 до $n$. Вагоны заезжают в тупик по одному в порядке номеров; в любой момент вагон, стоящий в тупике ближе всех к выезду, может выехать налево. Строка номеров вагонов в порядке выезда называется \term{стековой перестановкой}. Множество стековых перестановок длины $n$ обозначим $\mathcal S_n$.\\
Например, на рисунке $n=4$: вагон 2 уже выехал, 1 и 3 в тупике; если выедут 3 и 1, а затем 4, получится $2314$.}}

\begin{samepage}
\blok{\nomer{5}{\dagger}\hspace{\zazor}Пусть $\mathcal P_n$~— множество путей Дика в $(2n,0)$. Для каждого из множеств $\mathcal P_n$, $\mathcal D_n$, $\mathcal T_n$, $\mathcal S_n$:}
\punkt{а}{}{выпишите (нарисуйте) все его элементы при $n\leqslant 4$;}
\punkt{б}{}{выразите число его элементов при $n+1$ через числа элементов при меньших $n$;}
\punkt{в}{}{постройте биекции между ним и каждым из остальных.}
\end{samepage}\par\vspace{9pt}

\zad{6}{}{При каких $n$ число $C_n$ нечётно?}

\begin{samepage}
\blok{\nomer{7}{}\hspace{\zazor}«Бесконечный многочлен» $C(x)=C_0+C_1x+C_2x^2+\ldots$ назовём \term{производящей функцией} чисел~Каталана.}
\punkt{а}{}{Докажите, что $x\,C(x)^2=C(x)-1$.}
\punkt{б}{}{Найдите $C(x)$.}
\end{samepage}\par\vspace{9pt}

\newpage
\zvezdy

\tekst{Дальше~— несколько разных доказательств предположения из задачи 3в.}

\sboku{7.9cm}{@plohoj@}{\begin{samepage}
\blok{\nomer{8}{\dagger}\hspace{\zazor}Путь из $(0,0)$ назовём \term{плохим}, если он хотя бы раз опускается ниже оси (на рисунке~— плохой путь и прямая $y=-1$).}
\punkt{а}{}{Сколько плохих путей из $(0,0)$ в $(4,0)$? в $(6,0)$?}
\punkt{б}{}{Найдите $C_n$.}
\punkt{в}{}{Сколько путей Дика ведут в точку $(m,h)$?}
\punkt{г}{}{Найдите ответ в задаче 1.}
\end{samepage}\par}

\zad{9}{\circ}{С какой вероятностью на выборах из задачи 1 А ни разу не отставал от Б, если $a\geqslant b$?}

\sboku{3.9cm}{@krug@}{\begin{samepage}
\blok{\nomer{10}{}\hspace{\zazor}(\textit{Циклическая лемма.}) По кругу стоят $a$ чисел $+1$ и $b$ чисел $-1$, $a>b$. Начнём с одного из мест и по часовой стрелке прибавляем каждое число к сумме предыдущих.}
\punkt{а}{}{Сколько существует начальных мест, для которых все суммы положительны?}
\punkt{б}{}{Выведите отсюда формулу для $C_n$ другим способом.}
\end{samepage}\par}

\celoe{\begin{samepage}
\blok{\nomer{11}{}\hspace{\zazor}Звено пути назовём \term{верхним}, если хотя бы один его конец выше оси.}
\punkt{а}{}{Сколько путей из $(0,0)$ в $(6,0)$ имеют $0,2,4,6$ верхних звеньев?}
\punkt{б}{}{Сколько путей из $(0,0)$ в $(2n,0)$ имеют ровно $2k$ верхних звеньев?}
\punkt{в}{}{Выведите отсюда формулу для $C_n$.}
\end{samepage}\risunok{@shest@}}

\sboku{3.6cm}{@shestiug@}{\zad{12}{}{Постройте соответствие, которое каждому элементу $\mathcal T_{n+1}$ с отмеченной стороной, отличной от нижней, сопоставляет элемент $\mathcal T_n$ так, что каждый получается ровно $4n+2$ раза. Выведите отсюда формулу для~$C_n$ ещё одним~способом.}\par}

\zvezdy

\begin{samepage}
\blok{\nomer{13}{\circ}\hspace{\zazor}Монету бросают $2n$ раз. \term{Момент} $k$~— сразу после $k$-го броска, момент $0$~— начало. С какой вероятностью орлов и решек ни в какой момент, кроме $0$, не будет поровну?}
\punkt{а}{}{Найдите ответ при $2n=2,4,6,8$.}
\punkt{б}{}{Докажите, что таких исходов вдвое больше, чем путей Дика длины $2n-1$ (с любым концом).}
\punkt{в}{}{Найдите ответ при любом $n$.}
\end{samepage}\par\vspace{9pt}

\celoe{\begin{samepage}
\zadp{14}{\dagger}
\punkt{а}{}{Сколько путей из $(0,0)$ в $(6,0)$ имеют самую низкую точку на высоте $0,-1,-2,-3$?}
\punkt{б}{}{Постройте биекцию между путями длины $2n$, которые кончаются на оси, и путями Дика длины $2n$. Выведите отсюда ответ в задаче 13 другим способом.}
\end{samepage}\risunok{@nizh@}}

\begin{samepage}
\zadp{15}{}
\punkt{а}{}{Монету бросают 10\,000 раз. Докажите, что с вероятностью больше 0,99 орлов и решек после начала хотя бы раз будет поровну.}
\punkt{б}{}{Докажите, что ответ $u_n$ в задаче 13 удовлетворяет неравенствам \[\frac{1}{2\sqrt n}\leqslant u_n\leqslant\frac{1}{\sqrt{2n+1}}.\]}
\punkt{в}{}{С какой вероятностью орлов и решек впервые после начала станет поровну в момент $2n$? Докажите, что \[\frac{C_0}{4}+\frac{C_1}{4^2}+\ldots+\frac{C_{N-1}}{4^N}=\frac{1-u_N}{2}.\]}
\punkt{г}{}{Найдите $C(\frac14)$, где $C(x)$~— производящая функция из задачи 7. Как это согласуется с~пунктом~в)?}
\end{samepage}\par\vspace{9pt}

\zvezdy

\celoe{\begin{samepage}
\blok{\nomer{16}{\circ}\hspace{\zazor}Двое играют в орлянку: монету бросают несколько раз, орёл~— очко первому, решка~— второму; \term{путь игры} идёт звеном $\nearrow$ при орле и $\searrow$ при решке.}
\punkt{а}{}{При 4 и 6 бросках найдите вероятность того, что последний равный счёт был в~момент~$0,2,4,\ldots$}
\punkt{б}{}{Найдите вероятность того, что при $2n$ бросках последний равный счёт был в момент $2k$.}
\punkt{в}{}{Игра длится 100 бросков. Какова вероятность того, что последний момент равного счёта окажется среди моментов от 0 до 10? от 90 до 100? от 45 до 55?}
\end{samepage}\risunok{@posl@}}

\celoe{\sboku{7.6cm}{@arksinus@}{\begin{samepage}
\blok{\nomer{17}{\dagger}\hspace{\zazor}Игра длится $2n$ бросков. Будем говорить, что первый игрок \term{ведёт} во время броска, если соответствующее звено пути игры верхнее (см. задачу 11).}
\punkt{а}{}{Сколько путей длины 4, 6, 8 имеют $2k$ верхних звеньев?}
\punkt{б}{}{Постройте биекцию между путями длины $2n$ с $2k$ верхними звеньями и путями длины $2n$, у которых в последний раз счёт был равным в момент $2k$.}
\punkt{в}{}{Смоделируйте на компьютере 10\,000 игр по 1000 бросков. Постройте гистограмму доли времени, когда вёл первый, и нанесите на тот же чертёж график функции справа. В какой доле игр первый вёл больше 90\% времени?}
\end{samepage}\par}\risunok{@verh@}}

\celoe{\begin{samepage}
\blok{\nomer{18}{}\hspace{\zazor}\term{Сменой лидера} назовём момент, когда счёт равный, причём перед ним путь игры шёл по одну сторону от оси, а после~— по другую (на рисунке смены отмечены крестиками, равный счёт без смены~— кружками). Игра длится $2n$ бросков.}
\punkt{а}{}{Сколько путей длины 4, 6, 8 имеют $0,1,2,\ldots$ смен лидера?}
\punkt{б}{}{Докажите, что вероятность ровно $r$ смен равна $2\binom{2n-1}{n-1-r}\big/2^{2n-1}$. Какое число смен наиболее вероятно?}
\punkt{в}{}{Во сколько раз «лидер ни разу не сменился» вероятнее, чем «после начала счёт ни разу не~был~равным»?}
\end{samepage}\risunok{@smeny@}}
